% TOP_PROBLEM: {"id": 3403, "title": "One-way functions exist", "category_id": 15, "category": {"id": 15, "name": "computer_science", "display_name": "Computer Science", "description": "Computational complexity, algorithms, and theoretical CS.", "slug": "computer-science", "order_index": 15, "created_at": "2026-07-31T15:26:25.671Z"}, "status": "open"}
% FIRST_POSTED: null
% !TeX program = lualatex
% Compile twice: lualatex open_problems_500.tex
% All references and prose are embedded; no BibTeX database or external inputs.
\documentclass[11pt,a4paper]{article}
\usepackage[margin=24mm,headheight=14pt]{geometry}
\usepackage{fontspec}
\setmainfont{Latin Modern Roman}
\setsansfont{Latin Modern Sans}
\setmonofont{Latin Modern Mono}
\usepackage{amsmath,amssymb,mathtools}
\usepackage{unicode-math}
\setmathfont{Latin Modern Math}
\usepackage{microtype}
\usepackage{enumitem,xurl,needspace}
\usepackage[dvipsnames]{xcolor}
\usepackage{hyperref}
\hypersetup{unicode=true,colorlinks=true,linkcolor=MidnightBlue,urlcolor=MidnightBlue,
 pdftitle={Mathematical Research Notebook},pdfauthor={Source-annotated compilation},bookmarksnumbered=true}
\usepackage{bookmark}
\usepackage{tocloft}
\setlength{\cftsecnumwidth}{3em}
\cftsetpnumwidth{2.5em}
\cftsetrmarg{4em}
\usepackage{titlesec}
\titleformat{\section}{\Large\bfseries\raggedright}{\thesection}{0.7em}{}
\usepackage{fancyhdr}
\pagestyle{fancy}\fancyhf{}
\fancyhead[L]{\small\sffamily Open Problems Math | Research notebook}
\fancyhead[R]{\small 27 September 2026}
\fancyfoot[C]{\thepage}
\setlength{\parindent}{0pt}
\setlength{\parskip}{4pt plus 1pt}
\setlength{\emergencystretch}{3em}
\setcounter{tocdepth}{1}
\setcounter{secnumdepth}{1}
\setlist[itemize]{leftmargin=3em,itemsep=3pt,topsep=4pt,parsep=0pt}
\allowdisplaybreaks
\newcommand{\N}{\mathbb{N}}
\newcommand{\Z}{\mathbb{Z}}
\newcommand{\Q}{\mathbb{Q}}
\newcommand{\R}{\mathbb{R}}
\newcommand{\C}{\mathbb{C}}
\newcommand{\F}{\mathbb{F}}
\newcommand{\A}{\mathbb{A}}
\newcommand{\PP}{\mathbb P}
\newcommand{\E}{\mathbb{E}}
\newcommand{\eps}{\varepsilon}
\newcommand{\dd}{\,\mathrm d}
\newcommand{\sref}[2]{\hyperlink{source:#1:#2}{[S#2]}}
\newcommand{\eref}[2]{\hyperlink{extra:#1:#2}{[E#2]}}
% Turn the next switch off for a shorter reading copy. The quotations preserve
% every exactTarget field, including qualifications omitted from a short summary.
\newif\ifshowcatalogue
\showcataloguetrue
\newcommand{\cataloguescope}[1]{\ifshowcatalogue
 \paragraph{Exact scope in the supplied catalogue.}
 \begin{quote}\small #1\end{quote}\fi}
\usepackage{etoolbox}
\AtBeginEnvironment{itemize}{\small}
% Standalone entry; original section content is delimited for verification.
\begin{document}
\setcounter{section}{0}
%% BEGIN ORIGINAL SECTION CONTAINER
% ================================================================
% problemId: 3403
% Canonical problem ID: 3403
% exactTarget SHA-256: 9afb90e5fb1520e3a26583c1779c609dae28a39e36b697cf09190fc053e01c0a
\clearpage
\section*{\texorpdfstring{One-way functions exist}{One-way functions exist}}\label{3403}
\subsection{Definitions and mathematical statement}
A nonnegative function $\eta(n)$ is negligible if for every $c>0$ it is at most $n^{-c}$ for all sufficiently large $n$. The sought function family is uniformly computable in deterministic polynomial time and satisfies, for every probabilistic polynomial-time $A$,
\[
 \Pr_{x\leftarrow\{0,1\}^n,\,A}\!\bigl[A(1^n,f_n(x))\in f_n^{-1}(f_n(x))\bigr]
 \le \eta_A(n)
\]
for some negligible $\eta_A$. An output of the wrong input length is a failed inversion. The probability includes both a uniformly random input and the inverter's internal randomness. The inverter only has to find \emph{some} preimage; recovering the originally sampled input is unnecessary.
\par\smallskip\textit{Formulation and concept references:} \sref{3403}{1}.
\cataloguescope{Prove or disprove that there exists a deterministic polynomial-time computable function family f\_n:\{0,1\}\textasciicircum{}n->\{0,1\}\textasciicircum{}\{poly(n)\} such that every probabilistic polynomial-time algorithm inverts f\_n on a uniformly random input with only negligible probability.}
\subsection{Short English statement}
Is there an efficiently computable function whose outputs no efficient randomized algorithm can usually reverse, even when any preimage counts?
%% BEGIN RESEARCH INSERTION
\subsection{Research attempt}
\paragraph{Current frontier.}
\textit{Research check: 27 September 2026.}
The 2024 Hirahara--Lu--Oliveira manuscript gives characterizations of
one-way functions through public-coin time-bounded Kolmogorov complexity;
its Theorem 3 separately concerns \emph{infinitely-often} one-way
functions. Those quantifiers cannot be exchanged for the eventual
hardness required here. Liu--Pass (ITCS 2026, Theorems 4.1--4.2) gives a
boundary-hardness characterization for randomized time-bounded
complexity and a conditional characterization for its plain analogue.
These are equivalences or conditional reductions, not unconditional
constructions. \sref{3403}{2}, \sref{3403}{3}.
The attempt below instead tests an explicit polynomial map, distinguishing
its provable collision properties from its unproved inversion hardness.

\paragraph{Attack route.}
Three routes are natural: prove one of the meta-complexity hardness
conditions, convert worst-case NP-hardness into planted average-case
hardness, or analyze a concrete function directly. The first still needs
a complexity lower bound; the second must control the actual output
sampling distribution, not just arbitrary satisfiable instances.
We pursue the third with random quadratic systems whose coefficients
are part of the uniform input, not a separately generated secret key.

\paragraph{Attempt: a fully uniform candidate on every length.}
For $k\ge1$ let $m=3k$ and
\[
 q_k=1+k+\binom{k}{2},\qquad d_k=mq_k.
 \tag{8.1}
\]
A seed $Q\in\F_2^{d_k}$ describes $m$ independent multilinear quadratic
polynomials in $k$ variables, with all constant, linear and squarefree
quadratic coefficients included. For input length $N\ge d_1+1$, let $k$
be the largest integer with $d_k+k\le N$, and put $r=N-d_k-k$.
Parse the input as $(Q,x,u)$ with lengths $d_k,k,r$, and define
\[
 f_N(Q,x,u)=(Q,Q(x),u).
 \tag{8.2}
\]
For the finitely many smaller lengths use the identity map.
The output length is $N+2k$, evaluation is deterministic polynomial
time, and $k=\Theta(N^{1/3})$. Indeed, $d_k+k$ is a cubic polynomial
with positive leading coefficient, and maximality bounds $N$ above by
$d_{k+1}+k+1$. This is one algorithm on all lengths, with no advice.

For fixed $x\ne y$, the difference of one row's evaluations is a nonzero
linear functional of its coefficient vector: at least one linear
monomial distinguishes $x$ and $y$. Consequently
\[
 \Pr_Q[Q(x)=Q(y)]=2^{-3k},\qquad
 \Pr_Q[Q\text{ is not injective}]\le
 \binom{2^k}{2}2^{-3k}<2^{-k-1}.
 \tag{8.3}
\]
The second statement follows by a union bound over unordered pairs and
requires no independence between different pairs. Different seeds cannot
collide under (8.2), since the seed is copied to the output. For all but
an exponentially small fraction of seeds, the $x$-preimage is therefore
unique. This establishes a property of the proposed map, not one-wayness.

The necessary next statement is genuinely distributional: for every PPT
algorithm $A$, its probability of finding \emph{any} $x'$ with
$Q(x')=Q(x)$, given $(Q,Q(x))$ for uniform independent $Q,x$, would have
to be negligible in $k$. If this held, (8.2) would satisfy the target.
To justify the padding step, a proposed inverter at lengths $N$ can be
simulated at core length $k$ by sampling one of the
$O(k^2)$ lengths whose maximal parameter is $k$ and then sampling $u$.
Give each length probability at least the reciprocal of twice the
number of lengths: use a fixed block of random bits and map surplus
indices to the first length. This avoids unbounded rejection sampling.
A success $N^{-c}$ on infinitely many $N$ would give success at least
$k^{-3c-O(1)}$ on infinitely many $k$, contradicting the stated core
hardness. Thus this implication handles arbitrary length-sensitive
inverters, not merely an inverter that ignores padding. The core
hardness statement itself is not proved.

\paragraph{Stress test: injection is compatible with easy inversion.}
Replace the quadratic rows by independent affine rows $Q(x)=Ax+b$, with
$A\in\F_2^{3k\times k}$. The fixed-pair collision probability is
\emph{the same} $2^{-3k}$, so (8.3) still holds. Yet Gaussian elimination
finds a preimage in polynomial time on every output in the image, even
when $A$ is rank deficient. Thus both the near-injectivity bound and
its union-bound proof are incapable of distinguishing this easy family
from the quadratic candidate.

Nor is hardness on a uniform right-hand side the right hypothesis.
For each seed the image has at most $2^k$ elements in a space of size
$2^{3k}$, so an independent uniform $y$ lies in it with probability at
most $2^{-2k}$. Failure to solve $Q(x)=y$ under that sampling law would
usually mean that a preimage does not exist. In (8.2) it always exists.
Linearizing the monomials also fails to settle the issue: set
$z_{ij}=x_ix_j$, solve the resulting linear equations, and one must still
enforce all the relations $z_{ij}=x_ix_j$. Counting the many formal
linearized solutions neither proves these constraints hard nor supplies
an algorithm. In particular, no claim of cryptographic security is
inferred from degree two or from the number of equations.

\paragraph{A complementary obstruction to overcompression.}
For an arbitrary map $f$ on uniform inputs, put
$p_y=\Pr[f(U_N)=y]$. An inverter that samples $T$ independent inputs and
returns the first one matching a supplied output has failure probability
\[
 \sum_y p_y(1-p_y)^T\le \frac{|\operatorname{im}f|}{T+1}.
 \tag{8.4}
\]
For $0\le p\le1$, the maximum of $p(1-p)^T$ is at $p=1/(T+1)$ and is
at most $1/(T+1)$; summing proves the bound. A polynomial-size image is
therefore incompatible with one-wayness: choose a polynomial number of
trials larger than six times a polynomial upper bound for that image.
The candidate (8.2) has a much larger image, so this is an obstruction
to a different compression strategy, not its inverter.

Finally, $\mathsf P=\mathsf{NP}$ would invert every polynomial-time
family on its image by deciding successively whether a preimage extends
a chosen prefix. The queried relation is in NP and has polynomial-size
inputs. Establishing the missing hardness for (8.2) would thus already
imply $\mathsf P\ne\mathsf{NP}$. The converse is not available: sparse
worst-case difficulty need not be sampled with nonnegligible probability.

\paragraph{Outcome.}
\textbf{Outcome: Exploratory approach with unresolved gap.}
The proposed family meets the exact uniformity, input-length and
polynomial-output-length requirements, and its collision estimate and
all-length padding reduction are proved above. Its planted inversion
hardness has not been established. The affine countermodel shows why
the strongest easy estimate obtained here is not evidence of that
hardness. The map is a candidate to analyze, not a proved one-way
function or a claim of a new cryptographic primitive.
%% END RESEARCH INSERTION
\subsection{Sources}
\begin{itemize}\raggedright
\item[\textnormal{[S1]}]\hypertarget{source:3403:1}{} Shuichi Hirahara, Zhenjian Lu, and Igor C. Oliveira, One-Way Functions and pKt Complexity, IACR ePrint 2024/1388 (2024).
\url{https://eprint.iacr.org/2024/1388.pdf}.
{\small\textit{Catalogue formulation source.}}
\item[\textnormal{[S2]}]\hypertarget{source:3403:2}{} One-Way Functions and pKt Complexity — Hirahara, Lu and Oliveira.
\url{https://wrap.warwick.ac.uk/id/eprint/187935/1/WRAP-One-way-functions-pKt-complexity-24.pdf}.
{\small\textit{Catalogue research/status reference; not a proof endorsement.}}
\item[\textnormal{[S3]}]\hypertarget{source:3403:3}{} One-Way Functions and Boundary Hardness of Randomized Time-Bounded Kolmogorov Complexity — Liu and Pass.
\url{https://drops.dagstuhl.de/entities/document/10.4230/LIPIcs.ITCS.2026.97}.
{\small\textit{Catalogue research/status reference; not a proof endorsement.}}
\end{itemize}

%% END ORIGINAL SECTION CONTAINER
\end{document}
